\chapter{Introduction}

\section{Project Overview}
Jutsu Master is a real-time computer vision application that detects Naruto hand signs from live camera input, classifies gestures with a CoreML model, and maps valid sign sequences into interactive jutsu effects and gameplay states on iOS.

The full stack combines:
\begin{itemize}[leftmargin=2em]
\item dataset preparation and training in Python,
\item feature extraction with MediaPipe hand landmarks,
\item RandomForest classification exported to CoreML,
\item iOS runtime inference, sequence logic, visual effects, and audio.
\end{itemize}

\section{Problem Statement}
Hand-sign recognition in real-world camera conditions is difficult because of partial visibility, one-hand versus two-hand transitions, and pose variation between users. The project targets a robust unified model that can operate even when only one hand is visible.

\section{Technical Objectives}
\begin{itemize}[leftmargin=2em]
\item Build a single feature representation for one-hand and two-hand gestures.
\item Train a deployable classifier for 12 hand-sign classes.
\item Export and integrate the model into an iOS app pipeline.
\item Implement sequence-based jutsu detection and game logic.
\item Deliver a playable experience across free, tutorial, speed, and battle modes.
\end{itemize}

\section{Unified Feature Design}
Each hand contributes 21 landmarks with \(x,y,z\) coordinates.

\[
d = 2 \times 21 \times 3 = 126
\]

The model input vector is:
\[
\mathbf{x} = [\mathbf{h}_{\text{left}}, \mathbf{h}_{\text{right}}] \in \mathbb{R}^{126}
\]

If only one hand is detected, the missing side is filled with 63 zeros. This preserves a fixed input dimension while retaining one-hand usability.

\section{End-to-End Runtime Flow}
\begin{enumerate}[leftmargin=2em]
\item Capture a camera frame.
\item Detect hand landmarks (MediaPipe or Vision fallback).
\item Normalize and merge features into a 126-dimensional vector.
\item Run CoreML gesture inference.
\item Apply hold logic and sequence tracking.
\item Trigger jutsu effect, audio, and mode-specific game state updates.
\end{enumerate}

\placeholderfigure{system_pipeline}{figures/screenshots/system_pipeline.png}{System pipeline placeholder. Replace with a final architecture diagram or app pipeline screenshot.}

\section{Document Scope}
This report documents dataset preparation, model training, iOS architecture, gameplay logic, measured results, and current limitations, with figure placeholders for assets and screenshots.
